\documentclass[11pt]{article}
\usepackage{graphicx} % Required for inserting images
\usepackage{amsmath}
\usepackage{amssymb}
\usepackage{multirow}
\usepackage{booktabs}
% \usepackage{array}
\usepackage{tabularx}
\usepackage{fontspec}
\setmainfont{Average-Regular}
% \usepackage{times}
% Load the setspace package
\usepackage{setspace}
% Using \doublespacing in the preamble 
% changes the text to double-line spacing
\doublespacing

\newtheorem{theorem}{Theorem}
\newtheorem{definition}{Definition}
\newtheorem{conjecture}{Conjecture}

\DeclareMathOperator{\im}{im}
\DeclareMathOperator{\vproj}{proj}

\title{MATH598: A Review of Evolutionary Graph Theory}
\author{Eoin O'Hearn}
\date{\today}

\begin{document}
Eoin O'Hearn

Theory of Mind: Foreword

14 December 2025	
 
My first introduction to these ideas that have garnered my full attention was in PHL212: Ancient Philosophy taught by Seth Bordner. Before I started the class, I held the Ancient Greek philosophers in high regard, though I didn't know anything about them. So, I was delighted to learn that they not only revered my choice of major, mathematics, but also saw them as being more than a description of the world. Plato thought that mathematics was not just a tool for modeling the natural world but was actually the natural world in its purest form. Because these ideas of triangles, shapes, and relationships were everlasting, then they were more real as well—realer than anything that could be perceived in the physical realm. These mathematical forms are a step between the physical reality and the Land of Forms, the truest reality that only contains the perfect and unique Ideals that underpin our reality. Plato's understanding of the universe placed extreme value on the rational disciplines, namely philosophers and mathematicians. Needless to say, I walked out of that class with a big head and newfound respect for Plato, but also a new perspective on my major that planted the seeds for my monomania with consciousness.

If we accept the current scientific understanding of our world and how we arose, the existence of consciousness becomes quite puzzling. At the sub-atomic level, observation collapses the wave equation, energy condenses and localizes to form quarks that themselves make up the fundamental particles - electrons, protons, and neutrons. Atoms collide under the extreme pressure of burning suns, fusing together to produce heavier and heavier elements, as well as photons, heat, and different forms of radiation. Eventually, bonds are produced and molecular compounds float about, but their complexity is restrained by the force of entropy. It is not until self-replicating compounds emerged that complexity is able to be stored and increased over generations of molecular configurations. And finally, the familiar face of evolution appears with the first strands of RNA. Amino acids combine to form proteins and lipids combine to form membranes, all to further RNA's selfish aims. Proto-cells form proto-sensors, eventually able to capture and store the energy of the sun, leading to themselves being captured and stored by other proto-cells.

The first eukaryotic cells emerge as natural factories that minimize energy expenditure and maximize energy obtainment. But selection is cruel and cold, and there is warmth in community. The cells that worked together found and exploited this advantage, developing methods of adhering to each other that would become fundamental for the formation of organic tissue. Tissue became organs that became systems and the complexity of life exploded. Organisms found their niches and specialized, opening the door for other life forms to arise and construct their own niches. Selection drives evolution, forcing organisms to either adapt or die out. Entire ecosystems have risen and fallen during the history of Earth. Extinction acts as a palate cleanser as it wipes the landscapes clean, leaving few organisms. Those left over become the seeds of the next epoch and the cycle renews.

But at what point did consciousness arise? When did cells go from sensing sunlight to feeling it?  Why is it that we can introspect on our own experiences? The laws of physics seem to be capable of explaining the entire history of our species, not just as animals in a fight for survival but as a set of extremely complex interactions between self-replicating molecules. But we are left with an intense intuition that we are living, not just surviving. In fact, it seems that consciousness is a prerequisite for understanding our world, including the laws of physics. It seems hard to imagine how we would be where we are without the ability to introspect, learn, and cognitively understand the world around us.

And so, we arrive where humanity has been since we first became aware of our awareness. Different disciplines have created different lines of investigation to answer the different peculiarities of the mind.

The philosophers have largely attempted to solve the puzzling nature of the interaction between the mind and body. It seems that consciousness is of a fundamentally different material than the body, but the two are still able to interact. This curiosity has led them in pursuit of determining the prima materia, the fundamental nature of our reality. Dualists have proposed that everything is either mental or physical, drawing a clean distinction between the two. Materialists have denied the existence of a mental realm, leaning on the success of the laws of physics to explain the world around us. Panpsychists posit that reality is fundamentally consciousness because of the inextricable relation between observation and quantum mechanics. If an observer is required for all physical phenomena, how can anything else be more fundamental than consciousness itself?

The psychologists have opted for a different track and have attempted to develop ways of analyzing the minds of others. The success of Freud's psychoanalysis dominated the field for the first half of the 20th century, but behaviorism arose to contrast the subjective nature of Freudian nature. Behaviorists attempted to reduce the mind to an input-output machine, with no regards for the internal, subjective experience of the observer. The later fields of functionalism and cognitive neuroscience have filled in some of the holes left by the behaviorists, but many remain.

The biologists have been concerned with methods of detecting consciousness. The growing evidence of agential animals constructing their own niches and experiencing REM sleep indicative of internal dream states, along with many other cognitive functions, caused the field to sign the New York Declaration on Animal Consciousness, confirming consciousness in mammals and birds. In 1973 Dobzhansky published his essay, “Nothing in Biology Makes sense Except in the Light of Evolution” and in a similar way it seems that any explanation for consciousness will have to account for the consciousness found across biodiversity.

As I prepare to graduate in the Spring of 2026, it is reassuring to have found a problem that will accompany me for the rest of my life. I have never been content with our society's push for specialization, and the interdisciplinary nature of the problem will allow me to pursue many different careers and fields of knowledge while still being unified under one umbrella. I know this to be the area I want to continue in because it has felt not that I chose this area, but that I have found it. I wouldn't say that I was aimless, but that I couldn't quite grasp what unified my interests. I explored mathematics, biology, computer science, and many different majors looking for a subject I could latch onto, but it wasn't until I let go of the disciplines that I started to focus on the actual problems at hand. My decision to join the NEW College, which at the time was explained loosely by a desire to take a wide array of classes, has given me the vocabulary to finally talk about the things that interest me as well as the confidence to live at the edge of the unknown, the edge of chaos. So, while I imagine the trajectory of my life after graduation, I find solace in the words of Thoreau and Crighton.
 
“'But,' says one, 'you do not mean that the students should go to work with their hands instead of their heads?' I do not mean that exactly, but I mean something which he might think a good deal like that; I mean that they should not play life, or study it merely, while the community supports them at this expensive game, but earnestly live it from beginning to end. How could youths better learn to live than by at once trying the experiment of living?”
-Walden, p. 94
 
“And you never will, because those things don't exist. No matter how seriously people take them,” Thorne said. “A hundred years from now, people will look back at us and laugh. They'll say, 'You know what people used to believe? They believed in photons and electrons. Can you imagine anything so silly?' They'll have a good laugh, because by then there will be newer and better fantasies.” Thorne shook his head. “And meanwhile, you feel the way the boat moves? That's the sea. That's real. You smell the salt in the air? You feel the sunlight on your skin? That's all real. You see all of us together? That's real. Life is wonderful. It's a gift to be alive, to see the sun and breathe the air. And there isn't really anything else. Now look at that compass, and tell me where south is. I want to go to Puerto Cortés. It's time for us all to go home.”
-The Lost World, p. 416
\end{document}
